\documentclass[letterpaper]{article}
\usepackage{aaai26} % AAAI 2027 will likely use same or similar style
\usepackage{times}
\usepackage{helvet}
\usepackage{courier}
\usepackage[hyphens]{url}
\usepackage{graphicx}
\usepackage{amsmath}
\usepackage{amssymb}
\usepackage{booktabs}
\usepackage{multirow}
\usepackage{subcaption}
\usepackage{hyperref}

\title{Early Features Are All You Need: How Entropy Bootstrapping\\Determines Agent Strength in Competitive Self-Play}
\author{
    Julian Quick \\
    % Affiliation \\
    % email
}

\begin{document}
\maketitle

\begin{abstract}
In competitive self-play, Soft Actor-Critic (SAC) dominates Proximal Policy Optimization (PPO) by a 95-to-2 margin across 616 agents and six training paradigms---a gap that reward shaping, opponent diversity, and hyperparameter optimization cannot close.
We trace this dominance to a brief high-entropy bootstrapping phase: auto-tuned entropy temperature crashes from its initial value to near-zero within the first 10\% of training, and this transient window determines final agent strength.
Transplant experiments reveal that the advantage is encoded in the \emph{actor's} first hidden layer, not the critic: freezing the actor after bootstrapping produces stronger agents than continued training (57.9\% vs 38.9\% combined strength), and pairing the bootstrapped actor with a randomly initialized critic recovers 51.6\%.
The qualitative mechanism---a universal entropy decay schedule followed by actor feature crystallization---transfers from 2D tag to 3D MuJoCo Ant Sumo, though quantitative optima are environment-specific.
These findings invert the primacy bias narrative: in competitive settings, early features are not a liability to be periodically reset, but the most valuable thing the agent learns.
\end{abstract}

% ============================================================
\section{Introduction}
\label{sec:intro}
% ============================================================

% Hook: SAC dominates PPO in competitive self-play
% What everyone expects: entropy enables sustained exploration
% What we find: brief bootstrapping phase -> actor features -> critic is replaceable
% Contrast with primacy bias (Nikishin 2022) and critic plasticity (Ma et al. ICLR 2024)
% Road map

\textbf{TODO: Write introduction.}

% ============================================================
\section{Related Work}
\label{sec:related}
% ============================================================

% SAC and entropy in RL (Haarnoja et al. 2018, Meta-SAC)
% Self-play methods (Silver et al., OpenAI Five, AlphaStar)
% Algorithm comparison in MARL (MAPPO - Yu et al. NeurIPS 2022)
% Representation learning in RL (Garcin et al. ICLR 2025)
% Primacy bias and plasticity (Nikishin et al. 2022, Ma et al. 2024, Lyle et al. 2024)
% Entropy scheduling (count-based, adaptive)

\textbf{TODO: Write related work.}

% ============================================================
\section{Environments and Evaluation}
\label{sec:env}
% ============================================================

% Tag environment (30x30, obstacles, vision rays, 9 layouts)
% Ant Sumo (MuJoCo, circular arena, push-out win condition)
% Cross-evaluation gauntlet methodology
% Why within-run metrics are misleading

\textbf{TODO: Write environments section.}

% ============================================================
\section{SAC Dominates Competitive Self-Play}
\label{sec:dominance}
% ============================================================

% 616-agent cross-method gauntlet: SAC 78-81%, PPO below 51%
% Null results: zoo training, reward shaping, HPO
% Figure 1: gauntlet heatmap

\textbf{TODO: Write dominance results.}

% ============================================================
\section{The Entropy Schedule}
\label{sec:entropy}
% ============================================================

% Alpha trajectory: crash to ~0.003 in 500K steps
% Universal across 8 reward presets
% Counterfactual: auto-tuned >> no entropy >> fixed 0.1
% Cross-domain: same shape in Ant Sumo
% Figure 2: alpha trajectories
% Figure 3: counterfactual bars

\textbf{TODO: Write entropy schedule section.}

% ============================================================
\section{The Mechanism: Actor Feature Conditioning}
\label{sec:mechanism}
% ============================================================

% Init-alpha determines strength via inverted-U (tag)
% Q-value correlation: best agent has worst critic at 500K
% Transplant experiments (5-seed): actor carries strength
% Freeze-actor beats continued training
% First hidden layer carries the features
% Connection to primacy bias (inverted)
% Figure 4: mechanism results
% Figure 5: init-alpha

\textbf{TODO: Write mechanism section.}

% ============================================================
\section{Cross-Domain Validation}
\label{sec:crossdomain}
% ============================================================

% Ant Sumo: entropy schedule transfers (same alpha ~0.004)
% Auto-tuning helps in both domains
% Specific init-alpha optimum is environment-specific
% This decomposition (universal mechanism + specific tuning) is itself useful

\textbf{TODO: Write cross-domain section.}

% ============================================================
\section{Corner Camping and Environment Geometry}
\label{sec:geometry}
% ============================================================

% 9 layouts, corner camping scales with obstacle-corner proximity
% Wall-midpoints: high wall use, zero corner camping
% Corner camping inversely correlated with cross-eval strength
% Figure 6: geometry

\textbf{TODO: Write geometry section (may move to supplementary).}

% ============================================================
\section{Discussion}
\label{sec:discussion}
% ============================================================

% Connection to plasticity literature: invert Nikishin, contradict Ma et al.
% Why continued training degrades features: distribution shift from arms race
% Practical implications: always auto-tune, don't transfer specific alpha values
% Limitations: 2 environments, 2-layer MLPs

\textbf{TODO: Write discussion.}

% ============================================================
\section{Conclusion}
\label{sec:conclusion}
% ============================================================

\textbf{TODO: Write conclusion.}

% ============================================================
\bibliography{refs}
\bibliographystyle{aaai}

\end{document}
